\section{Numerical validation and exploratory observations}\label{app:numerics}
The computations in this appendix test implementations and suggest analytical questions;
they are not used in the analytical derivation. Reported agreements refer to the sampled
configurations and stated grids, not to unrestricted error bounds.

% Moved from the stationary section: numerical validation, kept out of the main line.
%======================================================================
\subsection{Implementation comparisons at finite wavelength}\label{app:validation}\label{sec:stationary-validation}
%======================================================================
The analytical statements above do not rest on numerics. To test the assembled representation
away from the gradient expansion, \path{solution/numerics/compare_cubic_quadrupole.py} compares
\eqref{eq:combined} with an independent forward calculation
(\path{solution/numerics/forward_cubic.py}) at $\bar t_\star=1.2$, $\bar t_o=3$, $\bar t_I=0.5$,
$|k|\simeq1$, so that $|k|\bar t_o\approx3$.
\begin{itemize}
\item \emph{Forward side.} The fields are solved through third order from (N) and the Cartesian (T)
  rows with the initial row. The constraint rows are not used. The Fermi--Walker triad is
  transported with $\dd E/\dd t=-SE$, and the photon geodesic is integrated by multilinear Picard
  iteration. The emission and clock roots and the exact normalisation follow. Time integrals are
  spectral and sky averages use Gauss quadrature.
\item \emph{Assembled side.} $J_{111}$ is the repository's second-ray-free direct measurement,
  with its nested integrals evaluated by quadrature. $S_{\rm meas}$ is
  \eqref{eq:measurement-source}, and $J_{\rm lin}$ is \eqref{eq:combined-linear}.
\end{itemize}
For three triples the two sides of \eqref{eq:combined} agree to $\le3\times10^{-9}$:
a noncoplanar triple with unequal frequencies, a triple with an exactly zero pair momentum, and
one with $|k_b+k_c|=0.05$. The pieces themselves are of order one to ten. In addition:
\begin{itemize}
\item $J_{111}$ agrees with the forward first-field-only map to $\le3\times10^{-8}$;
\item $J_{\rm lin}$ in constraint rows agrees with the Cartesian third-order field to $4\times10^{-14}$;
\item the forward map reproduces the exact channel \eqref{eq:dilation-channel} to $4\times10^{-8}$;
\item the no-probe-ray Hessian alone matches the forward mixed response to $10^{-9}$.
\end{itemize}
Two controls are resolved: replacing the Fermi--Walker triad by the symmetric one changes the
generic-triple value by $9\times10^{-5}$, and dropping the clock changes it by $0.24$.

\paragraph{What survives the observer dilation.} With the exact observer-centred identity of
\S\ref{sec:obs-dilation}, $\mathcal K_3=\mathcal K^c_3-\sum_{\rm pairs}\mathfrak D\mathcal K_2
-\sum_a\mathfrak D^2\mathcal K_1$, the same code isolates the genuinely cubic observer-centred kernel.
For the noncoplanar triple it is $0.9\%$ of $\mathcal K_3$: $+0.063$ of $-6.759$. For the
configuration $(k,q,-q)$ that enters $C^{(13)}$ it is $33\%$: $-2.29$ of $-6.89$. In these two
samples the dilation of the linear and quadratic maps carries most of the cubic response,
though the centred remainder is substantial in the second. Prop.~\ref{prop:obs-dilation} converts that dilation
part into $\frac12\sigma^2D^2C^{(11)}$ and a total derivative. The shell values are in
the exploratory spectra below.

\paragraph{Gradient coefficients.}
The forward calculation reproduces the degree-two and degree-four coefficients of
\eqref{eq:stationary-four-gradient} to $1.3\times10^{-9}$ and $1.6\times10^{-7}$,
at $\iota=0.5$ and $1$. Their difference reproduces the observer-local term to $10^{-6}$
in the script \path{solution/numerics/check_kappa4_shell.py}.
This is supplementary to the exact curvature and stationary calculations. At $x_o=0.3$, the leading centred coefficient gives
$-0.194$ in units of $\mathcal K_1$, against the numerical value $-0.195$; higher gradients
remain in the latter.

\paragraph{Top-order adjoint and linear-asymptotic comparisons.}
In \path{solution/numerics/explore_adjoint.py}, the momentum-source form of
Prop.~\ref{prop:adjoint} agrees with the direct cubic Green response to a few parts in
$10^6$ at $x_o=1,3$. The equivalent contact form, which differentiates the Hamiltonian
source numerically, agrees to $10^{-3}$--$10^{-4}$; the full map comparison at $x_o=1.5$
is limited by its grid to $4\times10^{-4}$. At $x_o=1$, $\mu=0.9$, the contact form has
large cancelling homogeneous, lapse and Hamiltonian contributions, approximately
$14$, $-73$ and $58$ in units of $T_\ell$, with an unrounded sum of $-0.24$.
These are implementation diagnostics, not additional analytical cancellations.
For the linear subhorizon transfer at $\ell=2$, $R=3$, the expansion through $O(x_o^{-3})$
has relative differences $1.8\times10^{-2}$, $3.5\times10^{-4}$ and $7.6\times10^{-8}$
from the exact transfer at $x_o=40,120,360$, respectively.

\subsection{Higher-gradient projection fits}\label{app:projection-fits}
The exact centred-$22$ proportionality of \S\ref{sec:obs-dilation} is established through
spectral degree six. Numerical fits suggest that it does not extend to degree eight:
\path{solution/numerics/centred22_ratio.py} gives ratios $52.66$ at $(e,o)=(1.2,3)$ and
$52.29$ at $(0.8,3.5)$, stable to $\le0.015$ under a change of fit points. The lower-order
ratios $5/3$ and $28/3$ are reproduced to $3\times10^{-6}$.
With the fixed template $\TF(\hat k\hat q+\hat q\hat k)$,
\path{solution/numerics/half_dilation_projection.py} gives relative mismatches from the
$-\mathfrak D\delta\mathcal K_1/2$ relation of $1.3\times10^{-4}$ at $k_0=0.1$ and
$2.4\times10^{-3}$ at $0.2$, suggesting a higher-order departure.
Neither fit proves the next analytical coefficient or an all-orders failure. A degree-eight
spectrum calculation needs both the four-gradient square and the two-/six-gradient
interference; that calculation remains open.



\subsection{Exploratory finite-wavelength spectra}
\paragraph{What is computed.}
By Thm.~\ref{thm:logmap}, every part of the fourth-order spectrum except
$\frac13\operatorname{Re}C_\ell[\theta,L_3]$ needs only second-order quantities. Those parts
are $\frac14C^{(22)}_\ell$, $\mathcal N_\ell$, $\mathcal M_\ell$ and $\frac14C_\ell[L_2,L_2]$.
\path{solution/numerics/explore_fourth_order.py} evaluates them for $\ell=2$ on the shell ($k_0=1$,
$\sigma=1$), at fixed $\bar t_o/\bar t_\star=3$ and $\bar t_I=\bar t_\star/2$. The ingredients
are:
\begin{itemize}
\item the exact second-order sources of \code{second\_order.py}, inserted into the Green
  formulas of \S\ref{sec:fields} and integrated by quadrature;
\item the ray, emission, clock and normalisation formulas of \S\ref{sec:obs};
\item Gauss--Legendre quadrature in $\mu$ and on the sphere.
\end{itemize}

\paragraph{Validation.}
\begin{enumerate}
\item The first-order part reproduces the entire-function $T_2$ of Prop.~\ref{prop:Tl} to
  $10^{-6}$ for $x_o\le6$, and to $1.2\times10^{-4}$ at $x_o=12$ (ray grid $n_t=160\,x_o$).
\item $\frac14C_2[\theta^2,\theta^2]$ equals the first term of $\mathcal N_2$ to all printed
  digits.
\item The dilation Ward identity holds at finite $k$ ($x_o=1.5$) for the complete second-order
  map: $[\mathcal K_2(k,q)+\mathcal K_2(k,-q)]/2=-k\partial_k\mathcal K_1$, with residual
  $4.5\times10^{-3}$ at $q=0.04$ and $1.1\times10^{-3}$ at $q=0.02$, i.e.\ $O(q^2)$. As negative
  controls, dropping the ray terms or $A_2$ gives $O(1)$ failures. The clock and endpoint
  terms vanish in this limit and are not tested by it.
\item The clock shift agrees with a direct root of the exact clock equation to eight digits.
\item The moving-emission endpoint term agrees, to all printed digits, with the mixed
  $z_az_b$ coefficient obtained by solving the emission condition
  $\ln t_e+\nu(t_e,X(t_e))=\ln\bar t_\star$ directly for real two-mode fields.
\item $A_2'$ is computed analytically from the Green representation: the boundary terms of
  the cumulative integrals cancel, $\nu_2''$ comes from the lapse equation, and
  $c_{\rm Ham}'$ is differentiated symbolically. It agrees with a finite-difference
  derivative to $10^{-4}$.
\item Moving the early-time cutoff of the Green integrals from $10^{-3}\bar t_o$ to
  $2\times10^{-5}\bar t_o$ changes the results in the seventh digit.
\item The ray quadrature error scales as $n_t^{-2}$. On a coarse grid
  ($n_t=320$) it gives a $0.4\%$ error in $T_2$ at $x_o=12$. With $n_t=160\,x_o$, the first-order error is below
  $10^{-4}$ for $x_o\le12$. Grid refinement in $\mu$, the sphere and the ray changes a node
  value at $x_o\approx7.9$ only in the sixth digit.
\item As $x_o\to0$, $\frac14C_2^{(22)}/\sigma^2C_2^{(11)}$ approaches $\frac{17}3$.
\item At every node of $T_\ell$ below, the pipeline gives $C^{(11)}_\ell\sim10^{-12}$ and
  $\operatorname{Re}C_\ell[\theta,\theta L_2]\sim10^{-6}$. Both vanish exactly there
  (Prop.~\ref{prop:nodes}).
\end{enumerate}
The convergence tests were made at $x_o\le3$, at one node near $x_o=7.9$, and at $x_o=12$ on
both trajectories, where $24\to40$ $\mu$ nodes and a $1.5\times$ finer ray grid change
$\frac14C^{(22)}_2$ only in the sixth digit. At $x_o=16$ the $\mu$ quadrature needs 32 nodes, and
the $x_o=16$ rows use them (see the note below the tables).

\paragraph{Results.} Absolute values with $\sigma=1$ and $\ell=2$, along the fixed-ratio trajectory
$x_\star=x_o/3$, $x_I=x_o/6$. The first-order error is at most $1.2\times10^{-4}$.
\begin{center}
\begin{tabular}{rrrrrr}
\toprule
$x_o$ & $C_2^{(11)}$ & $\frac14C_2^{(22)}$ & $\frac14C_2[L_2,L_2]$ & $\mathcal M_2$ & $\mathcal N_2$\\\midrule
0.5 & 0.00243 & 0.00867 & 0.00867 & $1.9\times10^{-6}$ & $-6.0\times10^{-7}$\\
1   & 0.0258  & 0.0172  & 0.0177  & $-2.9\times10^{-4}$ & $-6.7\times10^{-5}$\\
1.5 & 0.0596  & 0.388   & 0.396   & $-0.0119$ & $-3.1\times10^{-4}$\\
2   & 0.0423  & 2.61    & 2.64    & $-0.0453$ & $9.5\times10^{-5}$\\
3   & 0.116   & 6.26    & 6.32    & $-0.320$  & $-0.0012$\\
4   & 0.958   & 8.32    & 8.65    & $-2.10$   & $-0.048$\\
6   & 0.0458  & 5.23    & 4.96    & $0.697$   & $0.126$\\
8   & 0.00059 & 6.80    & 4.65    & $1.26$    & $0.746$\\
12  & 0.00036 & 2.43    & 2.44    & $-0.057$  & $0.020$\\
16$^\dagger$ & 0.0533  & 1.29    & 1.33    & $0.267$   & $0.035$\\
\bottomrule
\end{tabular}
\end{center}
Fixed-emission trajectory, $x_\star=1$ and $x_I=\frac12$:
\begin{center}
\begin{tabular}{rrrrrr}
\toprule
$x_o$ & $C_2^{(11)}$ & $\frac14C_2^{(22)}$ & $\frac14C_2[L_2,L_2]$ & $\mathcal M_2$ & $\mathcal N_2$\\\midrule
2  & 0.0340 & 2.12 & 2.14 & $-0.032$ & $-6.7\times10^{-5}$\\
4  & 0.946  & 8.31 & 8.53 & $-1.98$  & $-0.028$\\
8  & 0.0014 & 4.21 & 2.57 & $0.907$  & $0.563$\\
12 & 0.0302 & 1.76 & 1.63 & $-0.304$ & $0.041$\\
16$^\dagger$ & 0.00091 & 0.787 & 0.719 & $-0.061$ & $0.047$\\
\bottomrule
\end{tabular}
\end{center}
\noindent$^\dagger$At $x_o=16$, 32 $\mu$ nodes are needed. With 24 nodes, which suffice at $x_o\le12$,
$\frac14C_2^{(22)}$ comes out 3\% low (fixed ratio: 1.2514 against 1.28688) and 15\% low (fixed emission: 0.66617
against 0.787269). Increasing to 48 nodes leaves both values unchanged to six digits, and a $1.5\times$ finer ray grid has
no effect. The relative-momentum integrand therefore acquires structure between $x_o=12$ and $16$ that 24 nodes do not
resolve. The $x_o=16$ rows above use 32 nodes.

\emph{Fourth-order spectra at linear-transfer nodes} (Prop.~\ref{prop:nodes}; exploratory values, not analytical evaluations). At
these points $C_\ell=\eps^4\sigma^4\,c_\ell$ exactly through fourth order, with:
\begin{center}
\begin{tabular}{llll}
\toprule
 & trajectory & nodes $x_o$ of $T_\ell$ & $c_\ell=\frac14C^{(22)}_\ell/\sigma^4$\\\midrule
$\ell=2$ & $x_\star=x_o/3$ & 2.49867, 6.38035, 7.90609 & 5.686, 3.982, 7.285\\
$\ell=2$ & $x_\star=1$     & 2.51338, 6.33088, 7.50682 & 5.538, 2.732, 6.222\\
$\ell=3$ & $x_\star=x_o/3$ & 3.05037, 7.00073, 8.69752 & 3.886, 4.334, 2.353\\
$\ell=3$ & $x_\star=1$     & 3.05208, 6.92329, 8.48546 & 3.920, 2.582, 1.156\\
\bottomrule
\end{tabular}
\end{center}
The node positions are roots of the entire-function $T_\ell$. Further nodes for $x_o\le14$
are at 9.46677 and 11.96208 ($\ell=2$, fixed ratio), at 8.52178, 10.89081 and 13.06476
($\ell=2$, $x_\star=1$), at 10.60041 and 13.04624 ($\ell=3$, fixed ratio), and at 9.38590,
11.66731 and 13.83506 ($\ell=3$, $x_\star=1$).

\paragraph{Observations.} The following are leads for analysis, not conclusions.
\begin{enumerate}
\item \emph{Linear dominance, not perturbative breakdown.} Inside the horizon,
  $\frac14C_2^{(22)}$ stays $O(1)\sigma^4$ while $C_2^{(11)}$ is smaller and passes through
  nodes of $T_2$. With $r_{22}=C_2^{(22)}/(4\sigma^2C_2^{(11)})$, the quadratic contribution is
  smaller than the linear one only if $\eps^2\sigma^2\,r_{22}\ll1$, and $r_{22}\gg1$ for
  $x_o\gtrsim2$. Away from nodes the $13$ term also enters the total correction. At a node, the
  first nonzero prediction is
  the fourth-order one (Prop.~\ref{prop:nodes}). This is a statement about which order
  dominates a given multipole. It does not establish breakdown of the expansion, which would
  need the next order. Away from nodes, the $13$ term must also be included; its evaluation below changes the
  fourth-order correction materially.
\item \emph{Measurement nonlinearity.} $\mathcal N_2+\mathcal M_2$ is below about $5\%$ of
  $\frac14C_2[L_2,L_2]$ for $x_o\le3$, and 25--45\% at $x_o=4$--8.
\item \emph{Large cancelling pieces.} At small $x_o$ almost all of $\frac14C_2[L_2,L_2]$
  comes from the $A_2$ part. At larger $x_o$ the local part, the $A_2$ part and their
  interference are each large and strongly anticorrelated: at $x_o=12$ (fixed ratio) they
  are 12.5, 17.4 and $-27.5$, summing to 2.4. This suggests recombining
  the amplitudes \emph{before} squaring. For example,
  $\int A_{2,nn}'e^{iK\cdot\nh(\bar t_o-t)}\dd t=\big[A_{2,nn}e^{iK\cdot\nh(\bar t_o-t)}\big]_{\bar t_\star}^{\bar t_o}+iK\cdot\nh\int A_{2,nn}e^{iK\cdot\nh(\bar t_o-t)}\dd t$
  turns the field term into endpoint pieces. Those could combine with the observer-frame, ray
  and emission terms, before the separately squared norms are formed.
\item \emph{Resonances.} The $\mu$-integrand is smooth at these finite times, as the
  entire-function representation guarantees. It peaks at $K\approx0.8$--$1.0\,k_0$ in the
  tested configurations. Neither fact settles the asymptotics. A resonant layer of width
  $\propto1/t$ is regular at every finite time yet can dominate at late times, so a matched
  asymptotic analysis in angle and time is needed before any late-time claim. In these runs,
  $\frac14C^{(22)}$ divided by $T_2(K)^2$ is not universal in $x_o$.
\item \emph{Two different limits.} The main table follows $x_\star=x_o/3$ and
  $x_I=x_o/6$: increasing $k_0$ at fixed physical labels, which is the $k_0$-dependence the
  problem asks for. The late-time limit of \S\ref{sec:limits} instead keeps $x_\star$ and
  $x_I$ fixed. The second table follows that trajectory with $x_\star=1$, $x_I=\frac12$. The
  two must not be conflated.
\end{enumerate}
\paragraph{The cubic term and the complete fourth-order spectrum.}
\path{solution/numerics/explore_third_order.py} evaluates the remaining piece,
$\frac13\operatorname{Re}C_\ell[\theta,L_3]$, at $(k,q,-q)$. Its ingredients are:
\begin{itemize}
\item the exact third-order sources of \code{third\_order.py}, together with the
  third-order Hamiltonian source, extracted for this purpose;
\item the numerical second-order pair fields, with the exact $K=0$ kernel for the pair
  $(q,-q)$;
\item the Green formulas of \S\ref{sec:fields} for $A_3$.
\end{itemize}
The map is not taken from the hand-expanded formulas. It is solved from its definitions by a
numerical multilinear engine: the ray equation, the fields at the displaced ray, the
emission root, the clock equation, the composition with the random observation time, and
the normalisation. The checks are:
\begin{enumerate}
\item The engine's second-order map reproduces the hand-formula evaluator above to
  $1.2\times10^{-6}$, and its first-order map to $3\times10^{-9}$. These are two independent
  implementations.
\item The third-order dilation Ward identity
  $\mathcal K_3(k,q,-q)\to(k\partial_k)^2\mathcal K_1$ holds at finite $k$ ($x_o=1.5$). The
  residual is $3.3\times10^{-2}$ at $q=0.08$ and $8.2\times10^{-3}$ at $q=0.04$, i.e.\ $O(q^2)$.
  Dropping $A_3$ gives a $46\%$ failure.
\item Keeping only the scalar part of $A_3$ gives the same $\frac13\operatorname{Re}C_2[\theta,L_3]$
  as keeping its vector and TT parts, to every printed digit. This confirms that the
  contracted cubic response is purely scalar.
\item As $x_o\to0$, $\frac13\operatorname{Re}C_2[\theta,L_3]/\sigma^2C_2^{(11)}$ is $3.720$ at
  $x_o=0.15$ and $2.895$ at $0.3$; removing the $O(x_o^2)$ term extrapolates to $3.995$,
  against the expected $4$ (\S\ref{sec:sh}).
\item At the node $x_o=2.49867$ of $T_2$ the cubic term is $-2\times10^{-5}$, consistent with
  its exact vanishing (Prop.~\ref{prop:nodes}).
\end{enumerate}
With $C^{(4)}_2\equiv\frac14C^{(22)}_2+\frac13\operatorname{Re}C^{(13)}_2$, assembled from
Thm.~\ref{thm:logmap} as
$\frac14C^{(22)}_2+\frac13\operatorname{Re}C_2[\theta,L_3]+\operatorname{Re}C_2[\theta,\theta L_2]-C_2^{(11)2}/2\pi$,
the fixed-ratio trajectory gives ($\sigma=1$):
\begin{center}
\begin{tabular}{rrrrrr}
\toprule
$x_o$ & $C_2^{(11)}$ & $\frac14C_2^{(22)}$ & $\frac13\operatorname{Re}C_2[\theta,L_3]$ & $C^{(4)}_2$ & $C^{(4)}_2/C^{(11)}_2$\\\midrule
0.15 & $2.209\times10^{-5}$ & $1.204\times10^{-4}$ & $8.22\times10^{-5}$ & $2.03\times10^{-4}$ & 9.17\\
0.3  & $3.414\times10^{-4}$ & $1.653\times10^{-3}$ & $9.88\times10^{-4}$ & $2.64\times10^{-3}$ & 7.74\\
0.6  & 0.00475 & 0.0135 & $-0.00084$ & 0.0127 & 2.67\\
1    & 0.0258  & 0.0172 & $-0.159$   & $-0.141$ & $-5.5$\\
1.5  & 0.0596  & 0.387  & $-0.804$   & $-0.420$ & $-7.0$\\
2    & 0.0423  & 2.611  & $-0.581$   & 2.010    & 47.6\\
3    & 0.116   & 6.264  & $-7.57$    & $-1.57$  & $-13.5$\\
4    & 0.958   & 8.324  & $-38.7$    & $-32.2$  & $-33.6$\\
\bottomrule
\end{tabular}
\end{center}
The small-$x_o$ ratios approach $\frac{29}3=9.67$. Inside the horizon the $22$ and $13$
contributions are of the same size and partly cancel, and the net fourth-order correction to
the quadrupole changes sign several times. Away from nodes the cubic term is therefore
essential: the second-order part alone would give the wrong sign
at $x_o=1$--$1.5$ and at $3$--$4$. Convergence of the cubic term: at $x_o=3$ it is stable to five
digits under $16\to24$ $\mu$ nodes and under a finer sphere ($20\to28$ polar nodes) and ray
($300\to450$) grid. At $x_o=4$ it is stable to five digits under $16\to24$ $\mu$ nodes and to
four digits under a finer sphere ($24\to32$) and ray ($400\to600$) grid.

\paragraph{Independent recomputation of the cubic term.} The validated forward map of
\S\ref{sec:stationary-validation}, in \path{solution/numerics/shell_c13.py}, recomputes the
shell ratio
\[
 \frac{\frac13\operatorname{Re}C^{(13)}_2}{\sigma^2C^{(11)}_2}=\frac{\ens{\mathcal K_3(k,q,-q)}_{\hat q}}{\mathcal K_1(k)}
\]
on the fixed-ratio trajectory. Unlike the calculation above, it keeps all three sectors of $A_3$,
the Fermi--Walker frame, and fields solved without the Hamiltonian-algebraic scalar. It gives:
\begin{center}
\begin{tabular}{rrrrr}
\toprule
$x_o$ & $\mu$ nodes & cubic ratio & rounded-table ratio & centred part\\\midrule
0.3 & 8  & $2.8955$  & $2.89$ & $-0.195$\\
1   & 12 & $-6.156$  & $-6.13$ & $-2.96$\\
3   & 10 & $-67.29$  & $-67.5$ & ---\\
\bottomrule
\end{tabular}
\end{center}
The rounded-table column is $(C^{(4)}_2-\frac14C^{(22)}_2)/C^{(11)}_2$ computed from the
printed entries above. Allowing half a last printed digit in each input gives ratio intervals
$[2.8745,2.9076]$, $[-6.1651,-6.0986]$ and $[-67.8745,-67.1974]$, respectively.
The recomputed values lie within these rounding intervals; this is compatibility with the
older rounded table, not agreement to every displayed digit or a numerical error bound.
The newer $x_o=3$ result has not been separately refined in its $\mu$ quadrature.
Its linear quadrupole reproduces $C^{(11)}_2=\frac{4\pi}{25}\mathcal K_1^2$ ($0.0258$ and $0.116$).
The last column is
$\ens{\mathcal K^c_3}/\mathcal K_1$ from Prop.~\ref{prop:obs-dilation}; the rest is the dilation part. It grows
from $7\%$ of the cubic term at $x_o=0.3$ to about half at $x_o=1$. This is consistent with
$\mathcal K^c_3(k,q,-q)=0$ at leading gradient order and quantifies its breakdown inside the
horizon. As before, these are numbers, not analytical evaluations.

\subsection{Amplitude-level recombination and cancellation diagnostics}
To locate the cancellations, split the second-order log-map kernel into pieces $P_i$ and
measure $\kappa=\sum_i\|P_i\|^2/\|\sum_iP_i\|^2$. The norm is the $\ell=2$ projection integrated
over $\mu$. $\kappa\approx1$ means the pieces barely cancel; $\kappa\gg1$ means they cancel
heavily. The pieces are:
\begin{itemize}
\item the local shear products, which split into the $e^{2\zeta}$ term $2(H_a'+H_b')$ and
  the $H_aH_b$ term;
\item the $A_2'$ term;
\item the ray displacement $X\cdot\partial\sigma$ and the direction terms $\sigma(n,e)$;
\item the emission endpoint and the clock.
\end{itemize}
The results:
\begin{center}
\begin{tabular}{lrrrr}
\toprule
grouping & $x_o=1$ & 3 & 6 & 12\\\midrule
local $|$ $A_2$ $|$ ray $|$ endpoint $|$ clock & 11.8 & 1.01 & 12.6 & 21.8\\
$\tilde\gamma$ ($=$ local $+A_2$) $|$ ray $|$ endpoint $|$ clock & 6.3 & 0.91 & 4.9 & 9.9\\
$\tilde\gamma$ integrated by parts (observer, emission, bulk) $|$ ray $|\cdots$ & 6.6 & 6.0 & 43.7 & --\\
$A_2+HH+e^{2\zeta}+$direction $|$ displacement $|$ endpoint $|$ clock & 1.78 & 1.24 & 3.8 & 10.3\\
Eulerian: above field group $-\Psi\cdot\partial\sigma$ $|$ $(X+\Psi)\cdot\partial\sigma$ $|\cdots$ & 1.92 & 1.21 & 1.91 & 3.49\\
orthonormal: $A_2$ $|$ $-2\sigma(n,v)$ $|$ displacement $|$ endpoint $|$ clock & 1.73 & 0.99 & 6.2 & 13.9\\
orthonormal: $A_2+$displacement $|$ $-2\sigma(n,v)$ $|$ endpoint $|$ clock & 1.23 & 0.86 & 3.45 & 3.99\\
\bottomrule
\end{tabular}
\end{center}
Four observations follow.
\begin{enumerate}
\item The full second-order metric
  $\tilde\gamma=A_2+\frac12(H_aH_b+H_bH_a)+2(H_a+H_b)$ (so that $\sigma^{(2)}=\frac12\partial_t\tilde\gamma$)
  roughly halves the cancellation at every $x_o$. Part of it is an artefact of the
  exponential parametrisation $\gamma=e^{2\zeta}e^H$.
\item Grouping the $e^{2\zeta}$ shear term with the direction terms removes almost all
  cancellation near horizon crossing ($\kappa=1.8$ at $x_o=1$). In orthonormal directions
  $e^{\zeta}e$, the static conformal factor drops out of $\sigma(e,e)$. Prop.~\ref{prop:orthonormal} completes
  this analytically: in the fully orthonormal direction $e^{\zeta}e^{H/2}e$ the driver is
  $\frac12H'$ through third order on the shell. In the map the local products, the
  $e^{2\zeta}$ term and the $\sigma(n,e)$ terms combine exactly into the single piece
  $-H_b'(n,v_a)-H_a'(n,v_b)$; the numerics reproduce this identity to quadrature accuracy.
  That piece is small: its share of the norm is $0.13$, $0.04$, $0.76$ and $1.3$ at
  $x_o=1,3,6,12$, against about 5 for each of its constituents. Its three constituents were
  cancelling among themselves, not against the rest of the kernel. The orthonormal rows of the
  table show what is left. Near horizon crossing almost nothing cancels, and grouping $A_2$
  with the displacement gives $\kappa=1.23$ at $x_o=1$, the lowest of any grouping. The
  late-time cancellation between $A_2$ and the displacement is untouched: orthonormal
  directions do not address it (observation 4).
\item Integrating the field term by parts makes things worse inside the horizon: the observer
  endpoint and the bulk term cancel.
\item At late times the residual cancellation is between the field group and the photon
  displacement $X\cdot\partial\sigma$. The first-order kernel tends to a static profile,
  $Q\to6/k_0^2$ up to oscillations, so $H_1\to-\frac6{k_0^2}(\partial_i\partial_j-\frac13\delta_{ij}\nabla^2)g$.
  This is a pure static relabelling of the fluid labels, $X\to X-\frac3{k_0^2}\nabla g$. Its
  products appear both in $\tilde\gamma$ and in the ray displacement, and they must cancel in
  the observable. This suggests removing the static part of $H_1$ by an explicit
  relabelling before expanding. That would be the analogue of the cancellation of bulk
  displacements in large-scale-structure perturbation theory.
  A first test supports this. Write $\Psi_a=-\frac i2Q_a(t)\,k_ae^{ik_a\cdot X}$ for the
  first-order fluid displacement that removes the scalar part of $H_1$; it tends to
  $-\frac3{k_0^2}\nabla g$ at late times. Split the photon displacement term as
  $X\cdot\partial\sigma=(X+\Psi)\cdot\partial\sigma-\Psi\cdot\partial\sigma$, and group
  $-\Psi\cdot\partial\sigma$ with the field and direction terms. This lowers $\kappa$ from
  3.8 to 1.9 at $x_o=6$ and from 10.3 to 3.5 at $x_o=12$, and leaves it unchanged near horizon
  crossing (last row of the table). The test keeps only the first-order scalar displacement.
  Including the full first-order relabelling does not change this. That relabelling is
  $e=(1-\partial\Psi)(e^Y-v)$ with $v=\partial_t\Psi$, and it adds the Jacobian rotation of
  $\sigma$ and the fluid-velocity term. Kept as separate groups, the new Jacobian and
  direction pieces cancel each other (7.1 against 6.5 at $x_o=12$). Combined, they return the
  same $\kappa=3.5$. The residual is between the Eulerian field group and the photon's
  displacement relative to the fluid. A likely reason is that only space has been relabelled.
  In the Eulerian spatial coordinates the conformal factor becomes $\zeta_Y=\zeta(1-Qk^2/6)$,
  which still oscillates with $O(1)$ amplitude ($\to-g\cos x$) at late times. The
  temperature-defined time slicing is also displaced, because the radiation temperature
  oscillates with $O(1)$ amplitude inside the horizon. The natural next step is therefore a
  full change to longitudinal gauge, including the time shift, in which the potentials decay
  as $x^{-2}$ inside the horizon.
\end{enumerate}
The observer-centred decomposition $g=g(0)+[g-g(0)]$ was also tested. Exact dilation
covariance, $\Theta[\zeta+c]$ with momenta rescaled by $e^{-c}$, implies two identities. The
purely dilation parts of $\frac14C^{(22)}$ and $\frac13C^{(13)}$ combine exactly into
$\frac12\sigma^2D^2C^{(11)}$, with $D=k_0\,\dd/\dd k_0$ at fixed times. The dilation--centred
cross terms combine into a total derivative $-2D\operatorname{Re}X$, with
$X=\int\!\!\int PP\,A^c_2A_1^*$. This generalises the superhorizon $8=4+4$. At finite $k_0$,
however, $\frac12\sigma^2D^2C^{(11)}$ is not the dominant part of $C^{(4)}$: the remainder is
$-8.3\,C^{(11)}$ at $x_o=1$ and $-67\,C^{(11)}$ at $x_o=3$. The cross term of
Prop.~\ref{prop:obs-dilation} was then evaluated (\path{cross_term} in \path{solution/numerics/explore_fourth_order.py}), which
determines the centred part $C^{c,4}_2=C^{(4)}_2-\frac12\sigma^2D^2C^{(11)}_2+2D_{k_0}X_2$. In units of
$C^{(11)}_2$, fixed ratio:
\begin{center}
\begin{tabular}{rrrrr}
\toprule
$x_o$ & $C^{(4)}_2$ & dilation & cross & centred\\\midrule
0.15 & 9.17 & 7.89 & $-0.25$ & 1.53\\
0.3  & 7.74 & 7.54 & $-0.95$ & 1.15\\
0.6  & 2.67 & 6.16 & $-3.30$ & $-0.19$\\
1    & $-5.49$ & 2.79 & $-5.95$ & $-2.33$\\
1.5  & $-7.05$ & $-4.03$ & $-0.18$ & $-2.83$\\
2    & 47.6 & $-10.4$ & 28.0 & 30.0\\
3    & $-13.5$ & 53.9 & $-168.5$ & 101.1\\
\bottomrule
\end{tabular}
\end{center}
The small-$x_o$ entries follow the proposition: the centred part tends to $\frac53$, and the
cross term vanishes as $x_o^2$. Inside the horizon the three pieces are each large and
cancel, as at $x_o=3$, so the observer-centred split relocates the cancellation rather
than removing it. Its value is the exact superhorizon structure, not a simplification at
finite $k_0$.

The cubic term was evaluated by contraction, as follows. Only the
$q$-contracted third-order field is needed, $\bar A_3(k,t)=\int_qP(q)A_3(k,q,-q;t)$. The
third-order operator depends only on the total momentum $k$, and after the contraction $k$ is
the only remaining direction. Hence
$\bar A_{3,ij}=\bar h_3(k,t)(\hat k_i\hat k_j-\frac13\delta_{ij})$: the loop-averaged vector and
TT responses vanish. The new third-order work is therefore a scalar problem (lapse plus
scalar spatial part) with $q$-averaged sources. These sources, and the lower-order products in
$L_3$, still contain the second-order vectors and tensors, which are kept.

